% !TEX root = ../thesis.tex
% ==============================================================================
%  PERSIAN ABSTRACT & KEYWORDS (frontmatter/abstract.tex)
% ==============================================================================
%
%  This file contains the Persian abstract and keywords (چکیده و کلیدواژه‌ها).
%
% ------------------------------------------------------------------------------
%  HOW TO USE & CUSTOMIZE (راهنمای استفاده):
% ------------------------------------------------------------------------------
%  - Replace the text below with your thesis research abstract and keywords.
%  - The text below serves as a technical summary and feature overview of this template.
%  - Wrap your abstract text in \begin{persianabstract} ... \end{persianabstract}.
%  - Add your keywords using \thesiskeywords{کلیدواژه اول، کلیدواژه دوم، ...}.
%  - Spacing above the abstract can be adjusted in `styles/custom.tex`
%    via \setlength{\thesisabstractbeforespace}{...}.
% ==============================================================================

\begin{persianabstract}
نگارش پایان‌نامه‌ها و رساله‌های دانشگاهی به زبان فارسی نیازمند رعایت هم‌زمان شیوه‌نامه‌های نگارشی مصوب، تنظیم دقیق ساختار دوجهته در زی‌پرشین و پایداری فرآیند کامپایل است. این قالب، یک زیرساخت نرم‌افزاری مدرن و ماژولار برای تدوین پایان‌نامه‌های کارشناسی ارشد، رساله‌های دکتری و پروژه‌های کارشناسی بر اساس ضوابط دانشگاه صنعتی شریف ارائه می‌دهد. معماری سیستم بر مبنای تفکیک قطعی میان هسته‌ی پردازش حروف‌چینی (\code{styles/common.tex}) و لایه‌ی پیکربندی اختیاری کاربر (\code{styles/custom.tex}) بنا شده است تا شخصی‌سازی اجزای سند بدون تغییر در سازوکارهای پایه‌ای انجام پذیرد.

از قابلیت‌های فنی این قالب می‌توان به تمرکز فراداده‌ها در فایلی واحد و بازتولید هماهنگ آن در صفحات عنوان و فرم‌های اداری، سامانه‌ی خودکار تولید واژه‌نامه‌های دوزبانه مبتنی بر ترتیب الفبای یونیکد فارسی، اسکریپت‌های اختصاصی جهت صدور تک‌برگ فرم‌های تصویب‌نامه و اظهارنامه، و ماکروی درج مستقیم پی‌دی‌اف‌های امضاشده بدون تداخل حاشیه‌ها اشاره کرد. همچنین، ارتقای سبک کتاب‌شناسی به بسته‌ی طراحی شده‌ی \code{unsrtabbrv-modern.bst} امکان تولید خودکار پیوندهای کلیک‌پذیر شناسه‌ی شیء دیجیتال (\lr{DOI})، نشانی‌های وب و پیش‌چاپ‌های آرکایو را فراهم نموده است. یکپارچه‌سازی فونت‌های ادوبی کَسِلون، نیلوفر و شریف، استانداردسازی فواصل و علائم شماره‌گذاری، امکان تنظیم پانوشت‌های تک‌ستونی و دو ستونی، و گذر موفق از اعتبارسنجی‌های نحوی بدون هشدار کامپایلر، چارچوبی پایدار و استاندارد را برای پژوهشگران فراهم می‌سازد.

\thesiskeywords{قالب پایان‌نامه، دانشگاه صنعتی شریف، زی‌پرشین، حروف‌چینی دوزبانه، واژه‌نامه خودکار، مراجع تعاملی، خط‌لوله‌ی کامپایل.}
\end{persianabstract}
